\chapter{Исследовательская часть}

\section{Технические характеристики}

Технические характеристики устройства,на котором проводились замеры:

\begin{itemize}
	\item AMD Ryzen 7 5800U (16) @ 4.51 ГГц;
	\item Оперативная память: 16 Гб;
	\item Ядро ОС Linux 6.10.10-arch1-1.
\end{itemize}

При проведении замеров ноутбук был включен в сеть и были запущены только системные приложения и реализованные программы.

\section{Расчёт размера массива}

Расчёт размера массива по варианту производится по формуле:

\begin{equation}
	\label{variant}
	N = \frac{X}{8} + \begin{cases}
		X \% 1000,  &\text{если $\frac{X}{4} \% 10 == 0$,} \\
		\ (\frac{X}{4} \% 10) * (X \% 10) + (\frac{X}{2} \% 10), &\text{иначе},
	\end{cases}
\end{equation}
где $X = 8114$ – номер задачи в redmine.

Для моего варианта $N = 1053$

\section{Трудоёмкость алгоритмов}

В рамках данного исследования было проанализировано количество операций сравнения в циклах, зависимое от индекса искомого элемента в массиве. Результаты представлены в виде гистограмм: для линейного поиска — рисунок \ref{fig:lin-exp}, для бинарного поиска — рисунки \ref{fig:binid} и \ref{fig:bincmp}. На последнем графике столбцы упорядочены по высоте. Индекс -1 на всех графиках указывает на отсутствие искомого элемента в массиве.

\begin{figure}
	\centering
	\includegraphics[width=0.9\textwidth]{lin}
	\caption{Гистограмма зависимости числа сравнений от индекса элемента в линейном поиске}
	\label{fig:lin-exp}
\end{figure}

\begin{figure}
	\centering
	\includegraphics[width=0.9\textwidth]{binid}
	\caption{Гистограмма зависимости числа сравнений от индекса элемента в бинарном поиске}
	\label{fig:binid}
\end{figure}

\begin{figure}
	\centering
	\includegraphics[width=0.9\textwidth]{bincmp}
	\caption{Гистограмма зависимости числа сравнений от индекса элемента в бинарном поиске, отсортированная по количеству сравнений}
	\label{fig:bincmp}
\end{figure}

\pagebreak

\section{Вывод}

Был проведен анализ трудоемкости алгоритмов поиска в массиве.

Как видно из рисунков \ref{fig:lin-exp}–\ref{fig:bincmp}, линейный поиск значительно уступает бинарному по количеству операций сравнения. При размере массива в 1053 элементов максимальное количество сравнений для линейного поиска составляет 1053, тогда как для бинарного — всего 20.

Кроме того, как отмечается на рисунке \ref{fig:bincmp}, кривая трудоемкости бинарного поиска выглядит как логарифмическая прямая, что подтверждает его логарифмическую сложность O(logN).

\clearpage
